% Text set between the rows of an alignment, as in real papers: amsmath's
% \intertext and mathtools' \shortintertext put the display skips around the
% text as \vskip's of their own, and TeX sets no display skip of its own
% there. An alignment that opens a paragraph, and one that closes it.
\documentclass{article}
\usepackage{mathtools}
% one page with nothing but the text: checked against plain lualatex too (cases.ts)
\pagestyle{empty}
\begin{document}
The two maps are defined on the same domain, and their composition is
\begin{align}
  f(x) &= \sum_{k=0}^{n} a_k x^k + \int_0^x g(t)\,dt, \\
  g(x) &= \frac{d}{dx} \bigl(f(x) - a_0\bigr) - \sum_{k=1}^{n} k a_k x^{k-1}.
  \intertext{The second line follows from the first by the fundamental theorem of calculus, applied once to the integral term and term by term to the polynomial:}
  h(x) &= f(g(x)) + g(f(x)), \\
  h'(x) &= f'(g(x))\,g'(x) + g'(f(x))\,f'(x).
\end{align}
With the short form the text sits closer to the rows around it,
\begin{align*}
  u &= v + w, \\
  \shortintertext{so that}
  v &= u - w,
\end{align*}
and the paragraph goes on after the alignment, as it would after any display.

\noindent
\begin{align}
  p(t) &= q(t) + r(t), \\
  \intertext{and in the same way}
  q(t) &= p(t) - r(t)
\end{align}
opens this paragraph, and an alignment ends the next one:
\begin{align*}
  s &= \lim_{n\to\infty} s_n, \\
  \intertext{where the limit exists since the sequence is bounded and monotone, and}
  s_n &= \sum_{k=1}^{n} \frac{1}{k^2}.
\end{align*}

A paragraph after it, long enough to fill one line of the text width and then some more words.
\end{document}
